\scanpage{7}
\deleted{$X$ \unsure{loc.} de type fini sur}

\uline{Prop} Soit $X$ schéma \deleted{\unsure{localement|loc. noeth}} \deleted{\unsure{noethérien}} \inserted{régulier} sur corps $k$ \deleted{\unsure{de}} \uline{parfait}. Soit $k'$ une extension de $k$, supposons que \deleted{les ann} $X$ soit loc. de t. f. sur $k$ ou $k'$ ext. de type fini de $k$, de sorte que les anneaux locaux de $X_{k'}$ sont noeth. Alors $X_{k'}$ est régulier

$\Updownarrow$\quad \uline{Corollaire} Soit $A$ un anneau local régulier sur $k$, $k'$ une extension de $k$. On suppose que $k'$ \deleted{est \unsure{loc.}} \inserted{ou $A$} de t. f. \deleted{\unsure{sur}} $k$., \deleted{\unsure{ou|{}} $A$ est de t. f. sur $k$} Alors $A\otimes_k k'$ est régulier.

% margin note written sideways
\margin{\uline{NB} Ne supposons pas $X$ \unsure{ess.|loc.} de t. f. sur $k$, \deleted{Il reste vrai que} \inserted{Alors si} $X$ est régulier pour tt extension \inserted{finie} \deleted{de la base} corps de base, \deleted{alors} pour tt $Y$ ess. de t. f. sur $S$, $Y$ régulier $\Longrightarrow X\times_kY$ \uline{régulier}. On dit que $X$ est géom régulier.}
% margin note written sideways
\margin{Ex : ext. de corps \uline{\unsure{alg.|sép.}}}
\uline{Prop} $X$ schéma \inserted{ess loc de t. f.} sur $k$, $k$, ext: parfait. Conditions équivalentes

\diagram{\begin{tikzcd}[ampersand replacement=\&,row sep=.75em,column sep=.8em,execute at end picture={\draw ([xshift=-.75em]\tikzcdmatrixname-1-1.north west) .. controls +(0,-.55) and +(.24,0) .. ([xshift=-1.2em]\tikzcdmatrixname-2-1.west) .. controls +(.24,0) and +(0,.55) .. ([xshift=-.75em]\tikzcdmatrixname-4-1.south west);}]
\text{a)}\arrow[d,Rightarrow] \& {\forall\ \text{schéma }Y\text{ sur }k,\ Y\text{ régulier}\Longrightarrow X\times_kY\text{ régulier}}\\
\text{b)}\arrow[d,Rightarrow]\arrow[dd,bend right=37,Rightarrow] \& {\forall\ \text{ext. }k'/k,\ X_{k'}\text{ est régulier}}\\
\text{c)} \& {\deleted{\forall}\,X_{\unsure{k,|k'}}\text{ régulier}}\\
\text{d)} \& {\forall\ \text{ext. finie }\deleted{\unsure{k'}}\ \text{radicielle }k'\text{ de }k,\ X_{k'}\text{ est régulier.}}
\end{tikzcd}}

\uline{Dém} b) $\to$ a) car si $A\to B$ hom$^{\mathrm t}$ local \inserted{plat} d'anneaux loc. noeth, $A$ régulier et $B/mB$ régulier \inserted{$\Rightarrow B$ régulier}

c) $\Rightarrow$ b) car si $\overset{\mathfrak m}{A}\to\overset{\mathfrak n}{B}$ hom local plat d'anneaux loc, noeth, \unsure{ds|des} $B$ régulier $\Rightarrow A$ régulier \ldots

d) $\Rightarrow$ c) par \unsure{principe} de limite.

\uline{Déf} On dit alors que $X$ est géom. régulier [ou \uline{lisse} sur $k$ si $X$ est loc. de t. f. sur $k$]

$X\deleted{\unsure{{}_k|{}_x}}$ \deleted{géom régulier}\deleted{lisse}\inserted{géom régulier} $\Longrightarrow X$ régulier. La réciproque est vraie pour tt $X$ sur $k$ (ou seulement pour tt $X=\operatorname{Spec}k'$, $k'$ ext. finie radicielle de $k$) ssi $k$ est parfait.

$X$, \deleted{géom rég}\deleted{lisse}\inserted{géom rég} $/k\Longleftrightarrow X_{k'}$ géom régulier sur $k'$

$X$ et $Y$ \deleted{géom réguliers}\deleted{lisse}\inserted{géom rég} sur $k$ $\Longrightarrow X\times_kY$ géom. réguliers sur $k$.

$X$ \deleted{\unsure{lisse}} géom. régulier/$k$ en \uline{un point} $x\in X$ : $\operatorname{Loc}_x(X)$ est géom rég sur $k$ \inserted{ou $\underline O_{X,x}$ est géom rég}

On a les faits analogues aux précédents.
\margin{\uline{NB} Si $k(x)$ est séparable sur $k$ régulier $\Rightarrow$ géom. rég.}

\deleted{\uline{Théorème} L'ens des p$^{\mathrm{ts}}$ de $X$, où $X$ est géom. régulier est ouvert. (On est ramené au cas où \deleted{\unsure{$X$}} $k$ est parfait \unsure{(i.e.|(} $=$ \unsure{Alg.|rég.} clos).)}

\uline{Prop}\inserted{(Critère jacobien)} $X$ \inserted{loc de t f sur $k$} est \deleted{géom régulier}\inserted{lisse} en $x\in X$ ssi il existe un \inserted{entier $n\geq0$ et} morphisme $f:U_u\longrightarrow\mathbb E^n_k$ ($U$ vois. ouvert de $x$) tel que $f$ soit ``étale'' en $x$, i.e. [loc. en prés. finie] plat et à fibres étales

\uline{Corollaire} L'ens des p$^{\mathrm{ts}}$ où $X$ est \deleted{géom} lisse sur $k$ est ouvert.
% archive sheet number: 6

